% -----------------------------------------------------------------------------------
\section{The causal matrix and cross-problem axioms}\label{sec:matrix}

Theorem~\ref{thm:map15} shows that within-problem criteria yield only singleton and pairwise
obstructions. If a non-trivial impossibility in the style of Arrow exists, it must therefore come
from requirements that relate a theory's choices on \emph{different} problems: invariance,
independence, contraction. Such requirements need a common description of problems, which the
following reduced form provides.

\subsection{Reduced form}

\begin{definition}[Causal matrix]\label{def:matrix}
The \emph{causal matrix} of a problem $D$ is the square matrix indexed by $\Pol_D$,
\[
  M_D[\pi,\sigma]=\sum_{s\in S_D}\mu_\pi(s)\,U_D(s,\sigma),
\]
the value of running $\sigma$ while the environment responds as it would to $\pi$.
\end{definition}

Row $\pi$ is the causal evaluation of every policy under the beliefs that $\pi$ induces; column
$\sigma$ records how $\sigma$ fares under every such belief state; and the diagonal is the ex-ante
value, $M_D[\pi,\pi]=V_D(\pi)$. The matrix compresses the problem: it retains expectations under
each induced law, and it forgets the structural states. Two consequences should be kept in view. The
column-dominance axiom below compares expected values, whereas \cri{2} compares structural outcomes,
including states of probability zero; and \enquote{$\pi$ maximises its own row} is a global form of
ratifiability, which differs from the local, information-set form \cri{3} in multi-stage problems.
For problems with a single choice, the matrix and the structural criteria agree, and all the
witnesses below are of that kind.

\begin{theorem}[Realisability]\label{thm:realise}
Every finite real square matrix is the causal matrix of a fair, well-defined problem with a single
information set. Deleting a set of indices together with their rows and columns is realised by
deleting the corresponding actions while keeping the remaining laws.
\end{theorem}

\begin{proof}
Given $M\in\Rl^{I\times I}$, let the agent choose $i\in I$ at a single information set; let the states
be $S=I$; let an infallible predictor forecast the chosen index, $\mu_i=\delta_i$; and let
$U(s,j)=M_{sj}$. Then $M_D[i,j]=U(i,j)=M_{ij}$. The problem is fair and well defined, and deleting
actions leaves the laws of the others unchanged.
\end{proof}

In this realisation \cri{1} selects the maximal diagonal entries, and \cri{3} accepts exactly the
indices $i$ whose diagonal entry is maximal in row $i$. Rules on matrices therefore have an
unrestricted domain, as social welfare functions do in Arrow's theorem
\autocite{arrow1951}, and the framework contains every finite recommendation structure.

\subsection{Rules and axioms}

Fix an infinite universe $\mathbb L$ of labels. A \emph{matrix} is a pair $(I,M)$ with $I\subseteq\mathbb L$
finite and nonempty and $M\in\Rl^{I\times I}$; for $\varnothing\ne K\subseteq I$, $M|_K$ is the
\emph{principal restriction} to $K\times K$, which keeps labels. A bijection $\sigma\colon I\to I'$ acts by
$(\sigma\cdot M)_{\sigma(i)\sigma(j)}=M_{ij}$.

\begin{definition}[Rule]\label{def:rule}
A \emph{rule} assigns to every matrix a reflexive relation $\succeq_M$ on $I$, with strict part
$\succ_M$, and recommends its maximal elements,
\[
  G(M)=\{i\in I:\text{there is no }j\in I\text{ with }j\succ_M i\},
\]
which are required to be nonempty for every $M$.
\end{definition}

The convention makes the comparison axioms below constrain the same relation that determines
choice, and it prevents a rule that lacks an ordering from recommending nothing. A decision theory
satisfies \emph{reduction} (\ax{A0}) if its choice on $D$ depends only on the labelled matrix $M_D$,
so that it selects from $G(M_D)$ for some rule. A \emph{single-valued} deterministic selector
cannot be neutral: on the zero $2\times2$ matrix, swapping the indices fixes the input but moves
either possible answer. Choice correspondences avoid this separate obstruction.

\begin{definition}[Cross-problem axioms]\label{def:axioms}
A rule may satisfy the following conditions, for all matrices $M$ and all indices $i,j$.
\begin{description}[font=\normalfont\sffamily,leftmargin=2.4em,labelwidth=1.6em,labelsep=.8em,itemsep=.35em]
  \item[N] \emph{Neutrality.} For every bijection $\sigma\colon I\to I'$:
        $i\succeq_Mj\iff\sigma(i)\succeq_{\sigma\cdot M}\sigma(j)$.
  \item[F] \emph{Common affine invariance.} $\succeq_{aM+b\one\one^\top}=\succeq_M$ for $a>0$, $b\in\Rl$.
  \item[R] \emph{Row translations.} $\succeq_{M+b\one^\top}=\succeq_M$ for every $b\in\Rl^I$
        (adding $b_k$ to every entry of row $k$).
  \item[P] \emph{Column Pareto.} If $M_{k\sigma}\ge M_{k\tau}$ for every $k\in I$, strictly for some $k$,
        then $\sigma\succ_M\tau$.
  \item[L] \emph{Restricted Pareto.} The same, with $k$ ranging only over $\{\sigma,\tau\}$.
  \item[D] \emph{Diagonal optimality.} $G(M)\subseteq\argmax_iM_{ii}$.
  \item[J] \emph{Global ratifiability.} Every $i\in G(M)$ satisfies $M_{ii}\ge M_{ij}$ for all $j$.
  \item[B] \emph{Binary independence.} $i\succeq_Mj\iff i\succeq_{M|_{\{i,j\}}}j$.
  \item[S] \emph{Self-sovereignty.} If $M,M'\in\Rl^{I\times I}$ have the same rows $i$ and $j$, then
        $i\succeq_Mj\iff i\succeq_{M'}j$.
  \item[C] \emph{Contraction.} $G(M)\cap K\subseteq G(M|_K)$ for every nonempty $K\subseteq I$.
  \item[O] \emph{Ordering.} $\succeq_M$ is complete and transitive.
  \item[Q] \emph{Non-triviality.} For some $I$ there are $M,M'\in\Rl^{I\times I}$ with $G(M)\ne G(M')$.
\end{description}
\end{definition}

\ax{P} is the expectational counterpart of dominance; \ax{D} is \cri{1}; \ax{J} is ratifiability of
pure choices; \ax{S} is the condition of \textcite[§6]{briggs2010}, under which the comparison of two
options depends only on how each of them evaluates all options; \ax{C} is Chernoff's condition, Sen's
property~$\alpha$ \autocite{chernoff1954,sen1971}; \ax{B} is an independence of irrelevant alternatives
for binary comparisons, of the kind that \textcite{podgorski2022} imposes on overall-value theories,
required here across menu sizes. Neutrality is
stated for bijections between different label sets, because the proofs below transport choices from
one menu to another; invariances are imposed on the relation, not only on its maximal elements.
Ordering concerns the actual output relation at each input. Requiring only that each $G(M)$ be
rationalisable by some weak order would be vacuous, since any nonempty set can be placed in a top
indifference class; requiring one order to rationalise the choices from all restrictions of a
given matrix is stronger and already entails contraction \autocite{arrow1959}.

\begin{lemma}\label{lem:bimplies}
\ax{B} implies \ax{S} and \ax{C}.
\end{lemma}

\begin{proof}
If rows $i$ and $j$ coincide in $M$ and $M'$, so does the principal block on $\{i,j\}$, and \ax{B}
transfers the comparison. For \ax{C}, let $i\in G(M)\cap K$. Any $j\in K$ with $j\succ_{M|_K}i$ would
satisfy $j\succ_{M|_{\{i,j\}}}i$ and hence $j\succ_Mi$, contradicting the maximality of $i$ in $M$.
\end{proof}

The following monotonicity condition is used in Section~\ref{sec:exante}. Write $E_{ii}$ for the
matrix with a single entry $1$ in position $(i,i)$.

\begin{definition}[Diagonal monotonicity, \Mon]
If $M'=M+\delta E_{ii}$ with $\delta\ge0$, then $i\in G(M)$ implies $i\in G(M')$, and every true
comparison $i\succeq_Mj$ or $i\succ_Mj$ remains true in $M'$.
\end{definition}

\subsection{What row-translation invariance means}

The axiom \ax{R} will turn out to be decisive, and it is tempting to read it as the matrix form of
causal evaluation. It deserves a precise interpretation.

\begin{theorem}[Factorisation]\label{thm:factor}
A rule satisfies \ax{R} if and only if $\succeq_M$ depends on $M$ only through the row-normalised
matrix $Z(M)_{ij}=M_{ij}-M_{ii}$, equivalently through the within-row differences of $M$.
\end{theorem}

\begin{proof}
Row translations leave $Z$ unchanged. Conversely, if $Z(M)=Z(M')$, then
$M'_{ij}-M_{ij}=M'_{ii}-M_{ii}$ is constant along each row, so $M'$ is a row translation of $M$.
\end{proof}

On a restricted domain, \ax{R} is factorisation through the image of $Z$; the zero-diagonal
representative need not itself be admissible. With this, familiar criteria can be classified once a
matrix semantics is fixed for each.

\begin{center}
\small
\renewcommand{\arraystretch}{1.25}
\begin{tabularx}{\textwidth}{@{}>{\raggedright\arraybackslash}p{9.6em}
                                  >{\raggedright\arraybackslash}X
                                  >{\raggedright\arraybackslash}p{6.6em}@{}}
  \toprule
  criterion & matrix semantics & satisfies \ax{R}? \\
  \midrule
  fixed-prior causal decision theory & rank by $w^\top M$ for a fixed weight vector $w$ & yes \\
  pure global ratifiability & admit $i$ iff $M_{ii}\ge M_{ij}$ for all $j$ & yes, as an admissibility condition \\
  mixed ratifiability & $p\in\operatorname{BR}(p^\top M)$ & yes, as an equilibrium set \\
  deliberational causal decision theory & requires initial beliefs, update dynamics, and a final selection & depends on these \\
  evidential theory with reference-policy law & maximise $M_{ii}$, with $P(i,s)=\rho_i\mu_i(s)$ & no \\
  ex-ante policy maximisation & maximise $M_{ii}=V(i)$ & no \\
  \bottomrule
\end{tabularx}
\end{center}

Ratifiability's invariance does not establish that ratifiable options exist or supply a ranking.
For deliberational theories, fixed-belief causal comparisons with difference-based update dynamics
preserve row baselines, but arbitrary update dynamics need not; \textcite{joyce2012} distinguishes
equilibrium credences from the act subsequently chosen, and \textcite{arntzenius2008} identifies
the update rule as additional structure needed to define equilibrium. An \enquote{FDT-style} theory
is ex-ante maximisation under stipulated functional counterfactuals and so fails \ax{R} under that
stipulation \autocite{yudkowskysoares2017}.

\begin{proposition}\label{prop:regret}
\ax{R} does not characterise causal rationality: ranking by $r_i=M_{ii}-\max_jM_{ij}$ satisfies
\ax{N}, \ax{F}, \ax{R}, \ax{O}, yet on
\[
  M=\begin{pmatrix}0&10&-1\\10&0&-1\\0&0&-1\end{pmatrix}
\]
it chooses option $3$, which both other columns strictly dominate.
\end{proposition}

\begin{proof}
The scores are $(-10,-10,-1)$. Column $3$ is $(-1,-1,-1)^\top$, below both $(0,10,0)^\top$ and
$(10,0,0)^\top$ in every row.
\end{proof}

So \ax{R} means exactly invariance to row-specific baselines. It captures a feature shared by several
causal procedures, but it supplies neither causal expected-utility maximisation nor dominance.
